\documentclass[11pt,a4paper]{article}
\usepackage[T1]{fontenc}
\usepackage{lmodern}
\usepackage{amsmath,amssymb,amsthm}
\usepackage[margin=25mm]{geometry}
\usepackage{microtype}
\usepackage[hidelinks]{hyperref}
\newtheorem{theorem}{Theorem}
\newtheorem*{externaltheorem}{Faltings' theorem (invoked)}
\setlength{\parindent}{0pt}
\setlength{\parskip}{5pt}
\setlength{\emergencystretch}{2em}
\hypersetup{pdftitle={A finiteness result for a quartic Diophantine equation},
  pdfsubject={A rigorous correction of the proposed infinitude claim}}

\begin{document}
\begin{center}
  {\LARGE The equation $x^4+xy+y^3=1$}\\[5pt]
  {\large A finiteness result and the obstruction to infinitude}
\end{center}
\vspace{2pt}

\textbf{Answer.} The proposed infinitude statement is false. In fact, there
are only finitely many rational solutions, and therefore only finitely many
integer solutions.

\textbf{Scope.} The algebraic checks below are
proved explicitly. The final deduction invokes the genus formula for a smooth
plane curve and Faltings' theorem, both stated where used. Thus the conclusion
is unconditional, but this is \emph{not a fully self-contained proof of
finiteness}. A proof of the requested infinitude cannot be supplied because
that assertion is false. No claim of novelty is made.

\section*{1. Intuition: check the curve before seeking a family}

The dominant terms suggest the balance $y^3\approx -x^4$. This describes
the possible size of a large solution; it does not produce exact integer
solutions. Before attempting a recurrence or parametrization, the useful
structural question is whether the curve is singular. Singularities can
reduce the genus of a plane curve and sometimes permit a rational
parametrization. Here the degree is four, and all points turn out to be
smooth, so the genus is three. That is the decisive obstruction.

\section*{2. Explicit geometric checks}

Put $f(x,y)=x^4+xy+y^3-1$. Its projective closure $C$ is the curve
\begin{equation}
 F(X,Y,Z)=X^4+XYZ^2+Y^3Z-Z^4=0.
 \label{eq:projective}
\end{equation}
All smoothness and irreducibility checks are made over $\mathbb C$,
so they also cover nonreal points.

\textbf{The point at infinity is smooth.}
Setting $Z=0$ in \eqref{eq:projective} gives $X^4=0$. Hence the only point
at infinity is $P=[0:1:0]$. Since
\[
 F_Z=2XYZ+Y^3-4Z^3,\qquad F_Z(0,1,0)=1,
\]
$P$ is nonsingular.

\textbf{The curve is geometrically irreducible.}
Suppose that $F=GH$ over $\mathbb C$, where $G,H$ are homogeneous of
positive degrees $d$ and $4-d$. Restricting to $Z=0$ gives
\[
 G(X,Y,0)H(X,Y,0)=X^4.
\]
Neither restricted factor is zero. Since their product is a power of $X$,
homogeneity forces
\[
 G(X,Y,0)=cX^d,\qquad H(X,Y,0)=c^{-1}X^{4-d}
 \quad(c\ne0).
\]
Consequently $G(P)=H(P)=0$. The product rule would then give
$\nabla F(P)=G(P)\nabla H(P)+H(P)\nabla G(P)=0$, contradicting
$F_Z(P)=1$. Thus $F$ has no such factorization.

\newpage
\textbf{There are no affine singularities.}
An affine singularity would satisfy
\begin{equation}
 f=0,\qquad f_x=4x^3+y=0,\qquad f_y=x+3y^2=0.
 \label{eq:singular}
\end{equation}
The case $x=0$ forces $y=0$, but $f(0,0)=-1$, so $x\ne0$.
Substituting $y=-4x^3$ into the last equation yields
\[
 x+48x^6=0,\qquad x^5=-\frac1{48}.
\]
Also \eqref{eq:singular} gives $xy=-4x^4$ and
$y^3=-xy/3=4x^4/3$. Hence
\[
 0=f=x^4-4x^4+\frac43x^4-1=-\frac53x^4-1,
 \qquad x^4=-\frac35.
\]
Taking the ratio of the two nonzero powers forces
\[
 x=\frac{x^5}{x^4}=\frac{5}{144}.
\]
This is a positive real number and cannot have fourth power $-3/5$.
The contradiction excludes every affine singularity, including complex
ones. Together with the check at $P$, it proves that $C$ is smooth.

\section*{3. Why the solution set is finite}

We now use the standard \emph{genus formula}: a smooth, geometrically
irreducible projective plane curve of degree $d$ in characteristic zero has
genus $(d-1)(d-2)/2$. The checks above justify its application and give
\[
 g(C)=\frac{(4-1)(4-2)}2=3.
\]

\begin{externaltheorem}
If a smooth, projective, geometrically irreducible curve over a number
field $K$ has genus at least two, its set of $K$-rational points is finite.
\end{externaltheorem}

\begin{theorem}
The equation $x^4+xy+y^3-1=0$ has only finitely many rational solutions.
In particular, it has only finitely many integer solutions.
\end{theorem}
\begin{proof}
The curve $C$ is defined over $\mathbb Q$, is smooth and geometrically
irreducible, and has genus three. Faltings' theorem therefore makes
$C(\mathbb Q)$ finite. Each rational solution determines a point
$[x:y:1]\in C(\mathbb Q)$, and different pairs determine different points.
Thus the rational solution set is finite. Its subset of integer solutions
is finite as well.
\end{proof}

\section*{4. Explicit solutions and the limit of this argument}

Five integer solutions, verified by direct substitution, are
\[
 \boxed{(-1,-1),\quad(-1,0),\quad(-1,1),\quad(0,1),\quad(1,0).}
\]
They are exactly the solutions with $|x|\le1$: for $x=-1,0,1$, respectively,
the equation becomes $y^3-y=0$, $y^3=1$, and $y^3+y=0$.

The proof above does \emph{not} establish that these five points form the
complete integer solution set, and it supplies no explicit bound on the
coordinates. Neither is needed to refute infinitude. The precise result
proved here is finiteness, with its reliance on the two stated general
theorems made explicit.

\end{document}
